\chapter{Mengukur}\label{ch:g3:measure}

Berapa panjangnya, berapa beratnya, berapa isinya, pukul berapa
sekarang? Mengukur mengubah dunia menjadi bilangan --- asalkan semua
orang memakai satuan yang sama. Bab ini menyajikan satuan panjang,
\omterm{def:g2:measure:units}{massa} dan \omterm{def:g2:measure:units}{kapasitas}, lalu mengajarkan cara membaca jam.

\section{Panjang}

\begin{definition}[Satuan panjang]\label{def:g3:measure:length}
Satuan dasar panjang adalah \emph{meter} (m). Untuk panjang yang
kecil: sentimeter (cm) dan milimeter (mm); untuk jarak yang jauh:
kilometer (km).
\[
1 \text{ m} = 100 \text{ cm}, \qquad
1 \text{ cm} = 10 \text{ mm}, \qquad
1 \text{ km} = 1\,000 \text{ m}.
\]
\end{definition}

\begin{omfigure}
\begin{tikzpicture}[scale=1.3]
  % badan penggaris menjulur melewati tandanya agar tidak ada angka yang terpotong
  \draw[thick] (-0.35,0) rectangle (8.35,0.9);
  \foreach \i in {0,...,8} {
    \draw[thick] (\i,0.9) -- (\i,0.55);
    \node[font=\small] at (\i,0.32) {\i};
  }
  \foreach \i in {0,...,80} \draw (\i*0.1,0.9) -- (\i*0.1,0.72);
  \draw[omThm, very thick] (0,1.25) -- (3.6,1.25);
  \node[omThm, above, font=\small] at (1.8,1.3)
    {$3.6$ cm $= 3$ cm dan $6$ mm};
\end{tikzpicture}

{\small Membaca penggaris: tanda besar untuk sentimeter, tanda kecil
untuk milimeter. Bendanya mulai dari tanda $0$ --- bukan dari tepi
penggaris! --- dan berakhir di tanda kecil ke-$6$ setelah angka $3$.}
\end{omfigure}

\begin{example}[Memilih satuan yang tepat]\label{ex:g3:measure:choose}
Seekor semut: kira-kira $5$ mm. Sebatang pensil: kira-kira $18$ cm.
Sebuah pintu: kira-kira $2$ m. Berjalan ke kampung sebelah: kira-kira
$3$ km. Mengatakan ``panjang sebuah pensil adalah $18$'' tidak berarti
apa-apa --- satuannya adalah bagian dari jawaban.
\end{example}

\begin{method}[Mengubah satuan]\label{met:g3:measure:convert}
Untuk mengubah satuan, kalikan atau bagilah dengan $10$, $100$ atau $1\,000$:
\begin{enumerate}
  \item dari satuan besar ke satuan kecil, \emph{kalikan}:
        $3$ m $= 3 \times 100 = 300$ cm;
  \item dari satuan kecil ke satuan besar, \emph{bagilah}:
        $250$ cm $= 250 \div 100 = 2$ m dan $50$ cm;
  \item untuk menjumlahkan panjang, ubahlah dulu ke satuan yang
        \emph{sama}: $1$ m $+ 35$ cm $= 100 + 35 = 135$ cm.
\end{enumerate}
\end{method}

\section{Massa dan kapasitas}

\begin{definition}[Satuan massa dan kapasitas]\label{def:g3:measure:mass}
\omterm{def:g2:measure:units}{Massa} diukur dalam \emph{gram} (g) dan \emph{kilogram} (kg),
\omterm{def:g2:measure:units}{kapasitas} (berapa banyak isi sebuah wadah) dalam \emph{liter} (L)
dan \emph{sentiliter} (cL):
\[
1 \text{ kg} = 1\,000 \text{ g},
\qquad
1 \text{ L} = 100 \text{ cL}.
\]
\end{definition}

\begin{example}\label{ex:g3:measure:mass}
Sepucuk surat: kira-kira $20$ g. Sekantong tepung: $1$ kg. Sebuah
bak mandi: kira-kira $150$ L. Segelas air: kira-kira $20$ cL. Untuk
menyeimbangkan timbangan yang memuat kantong $1$ kg melawan $650$ g
apel, tambahkan $1\,000 - 650 = 350$ g di sisi apelnya.
\end{example}

\section{Membaca jam}

\begin{method}[Membaca waktu]\label{met:g3:measure:clock}
Pada jam berjarum:
\begin{enumerate}
  \item jarum \emph{pendek} memberi jamnya: bacalah angka yang
        paling akhir dilewatinya;
  \item jarum \emph{panjang} memberi menitnya: setiap tanda besar
        bernilai $5$ menit, jadi jarum panjang di angka $3$ berarti
        $15$ menit;
  \item setelah lewat setengah, kita juga dapat membacanya sebagai
        ``kurang dari jam berikutnya'': $10{:}50$ adalah ``pukul
        sebelas kurang sepuluh''.
\end{enumerate}
\end{method}

\begin{omfigure}
\begin{tikzpicture}[scale=1.0]
  % jam 1: 3:00
  \draw[thick] (0,0) circle (1.3);
  \foreach \h in {1,...,12} {
    \node[font=\footnotesize] at ({90-\h*30}:1.05) {\h};
    \draw ({90-\h*30}:1.22) -- ({90-\h*30}:1.3);
  }
  \draw[very thick, omDef] (0,0) -- (0:0.65);
  \draw[thick, omThm] (0,0) -- (90:1.05);
  \node[below, font=\small] at (0,-1.55) {3:00};
  % jam 2: 7:30
  \begin{scope}[xshift=4.4cm]
  \draw[thick] (0,0) circle (1.3);
  \foreach \h in {1,...,12} {
    \node[font=\footnotesize] at ({90-\h*30}:1.05) {\h};
    \draw ({90-\h*30}:1.22) -- ({90-\h*30}:1.3);
  }
  \draw[very thick, omDef] (0,0) -- ({90-7.5*30}:0.65);
  \draw[thick, omThm] (0,0) -- (270:1.05);
  \node[below, font=\small] at (0,-1.55) {7:30};
  \end{scope}
  % jam 3: 10:50
  \begin{scope}[xshift=8.8cm]
  \draw[thick] (0,0) circle (1.3);
  \foreach \h in {1,...,12} {
    \node[font=\footnotesize] at ({90-\h*30}:1.05) {\h};
    \draw ({90-\h*30}:1.22) -- ({90-\h*30}:1.3);
  }
  \draw[very thick, omDef] (0,0) -- ({90-10.83*30}:0.65);
  \draw[thick, omThm] (0,0) -- ({90-10*30}:1.05);
  \node[below, font=\small] at (0,-1.55) {10:50};
  \end{scope}
\end{tikzpicture}

{\small Tiga buah jam: jarum pendek biru untuk jam, jarum panjang
merah untuk menit. Pada 7:30 jarum jam berada \emph{di antara} 7 dan
8; pada 10:50 ia hampir mencapai 11.}
\end{omfigure}

\begin{example}[Lama waktu]\label{ex:g3:measure:duration}
Filmnya mulai pukul $14{:}40$ dan berakhir pukul $16{:}10$. Berapa
lama? Bilanglah dalam dua lompatan:
\begin{enumerate}
  \item dari $14{:}40$ ke $15{:}00$: $20$ menit;
  \item dari $15{:}00$ ke $16{:}10$: $1$ jam $10$ menit.
\end{enumerate}
Seluruhnya: $1$ jam $30$ menit. Hati-hati: satu jam berisi $60$
menit, jadi $16{:}10$ dikurangi $14{:}40$ tidak dapat dihitung
seperti $1610 - 1440$!
\end{example}

\section{Latihan}

\begin{exercise}[$\star$]\label{exo:g3:measure:1}
Satuan mana yang akan kamu pakai (mm, cm, m atau km) untuk: lebar
sebuah paku; tinggi teman sekelasmu; panjang halaman sekolah; jarak
antara dua kota?
\end{exercise}

\begin{exercise}[$\star$]\label{exo:g3:measure:2}
Ubahlah: $5$ m menjadi cm; \quad $70$ mm menjadi cm; \quad $4$ km
menjadi m; \quad $600$ cm menjadi m.
\end{exercise}

\begin{exercise}[$\star$]\label{exo:g3:measure:3}
Ukurlah dengan penggarismu (sampai milimeter terdekat) sisi sebuah
segitiga yang digambar teman sekelasmu, lalu hitunglah panjang
seluruh tepinya.
\end{exercise}

\begin{exercise}[$\star$]\label{exo:g3:measure:4}
Hitunglah, dengan mengubah satuannya dulu: $2$ m $+ 45$ cm (dalam cm);
\quad $1$ km $+ 250$ m (dalam m); \quad $3$ cm $+ 8$ mm (dalam mm).
\end{exercise}

\begin{exercise}[$\star$]\label{exo:g3:measure:5}
Ubahlah: $3$ kg menjadi g; \quad $2\,500$ g menjadi kg dan g; \quad
$4$ L menjadi cL; \quad $350$ cL menjadi L dan cL.
\end{exercise}

\begin{exercise}[$\star$]\label{exo:g3:measure:6}
Sebuah resep memerlukan $750$ g tepung. Mona punya sekantong $1$ kg.
Apakah tepungnya cukup? Berapa yang akan tersisa?
\end{exercise}

\begin{exercise}[$\star$]\label{exo:g3:measure:7}
Bacalah waktunya: jarum pendek antara $2$ dan $3$, jarum panjang di
angka $6$; jarum pendek di angka $9$, jarum panjang di angka $12$;
jarum pendek hampir di angka $5$, jarum panjang di angka $10$.
\end{exercise}

\begin{exercise}[$\star$]\label{exo:g3:measure:8}
Sekolah mulai pukul $8{:}30$ dan pelajaran pagi berakhir pukul
$11{:}45$. Berapa lama pagi itu? (Bilanglah dengan lompatan, seperti
pada \cref{ex:g3:measure:duration}.)
\end{exercise}

\begin{exercise}[$\star$]\label{exo:g3:measure:9}
Keretanya berangkat pukul $9{:}50$ dan perjalanannya memakan waktu
$45$ menit. Pukul berapa ia tiba? (Lompatlah ke $10{:}00$ dulu.)
\end{exercise}

\begin{exercise}[$\star\star$]\label{exo:g3:measure:10}
Sebuah botol berisi $1$ L. Pandu menuangkannya ke gelas berisi $20$
cL. Berapa gelas penuh yang dapat ia tuangkan? (Ubahlah satuannya dulu.)
\end{exercise}

\begin{exercise}[$\star\star$]\label{exo:g3:measure:11}
Lina berjalan $600$ m ke sekolah, lalu pulang, setiap hari. Berapa
meter yang ia tempuh dalam satu minggu sekolah yang $5$ hari?
Berikan jawabannya dalam meter, lalu dalam kilometer.

\end{exercise}
